\chapter{Idioms}
\label{chap:idioms}
\section{C++ Idioms}
An idiom can be considered as the most natural way to express something 
or to solve a problem. More specifically, an idiom is a group of code 
fragments sharing an equivalent semantic role, which recurs frequently in 
one or more programming languages or libraries.
\newline
A comprehensive list of C++ Idioms can be found at~\cite{wiki-idiom-list}.
See also~\cite{youtube-cpp-idioms-everyone}.

\subsection{PImpl Idiom}
See~\cite{cppreference-pimpl}.
\newline
"Pointer to implementation" or "pImpl" is a C++ programming technique 
that \textbf{removes implementation details of a class from its object representation} 
by placing them in a separate class, accessed \textbf{through an opaque pointer}.
\newline
This idiom is also known as \textbf{\texttt{Compiler Firewall Idiom}}.
See also~\cite{wiki-pimpl},~\cite{gotw-24-compilefirewalls},~\cite{gotw-100-compilfirewalls}.

\subsubsection*{Problematic Scenario}
Whenever a class definition changes in a header, all users of that class 
must be recompiled. The reason is the textual inclusion on which 
the C++ build model is based. The user is assumed to know two features 
which could be affected by private members: \textbf{size and layout} 
and \textbf{functions}.

\subsubsection*{Solution}
\begin{lstlisting}[language=C++]
// Pimpl idiom - basic idea
class widget {
    // :::
private:
    struct impl;        // things to be hidden go here
    impl* pimpl_;       // opaque pointer to forward-declared class
};    
\end{lstlisting}

%%%
\subsection{Fast PImpl Idiom}
\label{subsec:fast-pimpl}
See~\cite{gotw-28-pimpl}.
\newline
Classic PImpl pays for an extra heap allocation and an extra indirection.
The \textbf{Fast PImpl} idea is to keep the opaque implementation, but 
store it in a fixed-size buffer \textbf{inside} the object (or otherwise 
avoid the heap) so that construction stays cheap while the header still 
hides the real layout.
\newline
Trade-off: you must pick a buffer size/alignment large enough for 
\texttt{impl}, and you typically still need out-of-line destructor / special 
members because \texttt{impl} is incomplete in the header.

\begin{lstlisting}[language=C++]
#include <cstddef>
#include <new>

class Widget {
public:
    Widget();
    ~Widget();
    Widget(const Widget&) = delete;
    Widget& operator=(const Widget&) = delete;

    void draw();

private:
    struct Impl;
    static constexpr std::size_t buf_size = 64;
    alignas(std::max_align_t) unsigned char buf_[buf_size];

    Impl* pimpl() { return reinterpret_cast<Impl*>(buf_); }
    const Impl* pimpl() const { return reinterpret_cast<const Impl*>(buf_); }
};

// widget.cpp
struct Widget::Impl {
    int x = 0;
    void draw() { /* ... */ }
};

Widget::Widget() { static_assert(sizeof(Impl) <= buf_size); new (buf_) Impl{}; }
Widget::~Widget() { pimpl()->~Impl(); }
void Widget::draw() { pimpl()->draw(); }
\end{lstlisting}

%%%
\subsection{Handle/Object Idiom}
\label{subsec:handle-object}
See~\cite{eliens-handle-body-idiom}.
\newline
Separate a thin public \textbf{handle} (what clients use) from a heavy 
\textbf{body/object} (where the real state and work live). The handle 
usually holds a pointer/reference to the body and forwards operations.
\newline
This is closely related to PImpl / Bridge: clients depend on a stable 
handle interface; the body can change, be shared, or be replaced without 
rewriting client code.

\begin{lstlisting}[language=C++]
class StringBody;                 // heavy representation

class String {                    // handle
public:
    String(const char* s);
    String(const String& other);
    String& operator=(String other); // often with copy-and-swap
    ~String();

    char operator[](std::size_t i) const;
    void append(const char* s);

private:
    StringBody* body_;
};
\end{lstlisting}

\subsection{Copy and Swap Idiom}
\label{subsec:copy-and-swap}
See~\cite{stackof-copyandswap},~\cite{gotw-59-exceptsafeclassdesI}.
\newline
This idiom aims at creating a strong exception-safe implementation of 
the overloaded assignment operator. 
\newline
This idiom makes use of the \textbf{RAII} idiom in order to acquire 
the new resource before it forfeits its current resource. If the 
acquisition is successful, the old resource is released; if it throws, 
the object is left unchanged (\textbf{strong exception guarantee}).

\begin{lstlisting}[language=C++]
class Buffer {
public:
    Buffer(std::size_t n);
    Buffer(const Buffer& other);
    ~Buffer();

    friend void swap(Buffer& a, Buffer& b) noexcept {
        using std::swap;
        swap(a.data_, b.data_);
        swap(a.size_, b.size_);
    }

    Buffer& operator=(Buffer other) noexcept { // pass by value = copy/move
        swap(*this, other);
        return *this;
    }

private:
    char* data_;
    std::size_t size_;
};
\end{lstlisting}

\subsection{RAII - Resource Acquisition is Initialization}
\label{subsec:raii}
See~\cite{wiki-raii}.
\newline
\textbf{RAII} ties a resource's lifetime to an object's lifetime: 
acquire in the constructor, release in the destructor. Stack unwinding 
then releases resources automatically on both success and exception paths.
\newline
Typical uses: locks, file handles, memory, sockets. Standard library 
examples include \texttt{std::unique\char`_ptr}, \texttt{std::lock\char`_guard}, 
\texttt{std::fstream}.

\begin{lstlisting}[language=C++]
class File {
public:
    explicit File(const char* path) : f_{std::fopen(path, "rb")} {
        if (!f_) throw std::runtime_error("open failed");
    }
    ~File() { if (f_) std::fclose(f_); }

    File(const File&) = delete;
    File& operator=(const File&) = delete;

    std::size_t read(void* buf, std::size_t n) {
        return std::fread(buf, 1, n, f_);
    }

private:
    std::FILE* f_;
};

void use() {
    File f("data.bin");          // acquire
    // ...                        // always released when f goes out of scope
}
\end{lstlisting}

\subsection{DBDII - Dynamic Binding During Initialization Idiom}
\label{subsec:dbdii}
See~\cite{isocpp-calling-virtuals-from-ctor-idiom}.
\newline
Calling a \texttt{virtual} function from a constructor/destructor does 
\textbf{not} dispatch to the most-derived override: the dynamic type is 
still under construction. DBDII avoids that trap by finishing construction 
first, then invoking a virtual ``post-constructor'' (often a factory 
\texttt{create()} that returns a smart pointer and then calls 
\texttt{init()}).

\begin{lstlisting}[language=C++]
class Base {
protected:
    Base() = default;
public:
    virtual ~Base() = default;
    virtual void post_construct() = 0;   // safe: object fully constructed

    template <class T, class... Args>
    static std::unique_ptr<T> create(Args&&... args) {
        auto p = std::unique_ptr<T>(new T(std::forward<Args>(args)...));
        p->post_construct();             // dynamic binding OK here
        return p;
    }
};

class Derived : public Base {
    friend class Base;
    Derived() = default;
public:
    void post_construct() override { /* use virtuals safely */ }
};

auto d = Base::create<Derived>();
\end{lstlisting}

\subsection{Named Constructor Idiom}
\label{subsec:named-ctor-idiom}
See~\cite{isocpp-named-ctor-idiom}.
\newline
Make constructors private/protected and expose clearer static factory 
functions with meaningful names. This documents intent, can enforce 
invariants, and can return by value / smart pointer without overloading 
ambiguity.

\begin{lstlisting}[language=C++]
class Point {
public:
    static Point cartesian(double x, double y) { return Point{x, y}; }
    static Point polar(double r, double theta) {
        return Point{r * std::cos(theta), r * std::sin(theta)};
    }

private:
    Point(double x, double y) : x_{x}, y_{y} {}
    double x_, y_;
};

auto p = Point::polar(1.0, 3.14159 / 4);
\end{lstlisting}

\subsection{Construct On First Use Idiom}
\label{subsec:construct-on-first-use}
See~\cite{isocpp-const-first-use-idiom}.
\newline
Avoid the static initialization order fiasco: instead of a namespace-scope 
object whose constructor may run before another translation unit is ready, 
wrap the object in a function-local \texttt{static} so it is constructed 
on first call (and destroyed at program exit in a well-defined way for 
that object).

\begin{lstlisting}[language=C++]
// Bad across TUs: relative init order of a and b is undefined
// extern Log log;
// Log& get_log() { return log; }

Log& get_log() {
    static Log log;   // constructed on first call
    return log;
}

void f() { get_log().write("hi"); }
\end{lstlisting}

\subsection{Nifty Counter Idiom}
\label{subsec:nifty-counter}
See~\cite{isocpp-nifty-counter-idiom},~\cite{wikibook-nifty-counter-idiom}.
\newline
Used by the standard streams (\texttt{std::cout} \& co.): each translation 
unit that includes the iostream header gets a static object whose 
constructor/destructor bump a shared counter. The first construction 
initializes the stream objects; the last destruction tears them down. 
That way streams are usable during dynamic initialization of other 
statics in TUs that included the header.

\begin{lstlisting}[language=C++]
// Simplified sketch of the idea (not the real library code)
class IostreamInit {
public:
    IostreamInit() {
        if (count_++ == 0) { /* construct cin/cout/cerr */ }
    }
    ~IostreamInit() {
        if (--count_ == 0) { /* destroy cin/cout/cerr */ }
    }
private:
    static int count_;
};

static IostreamInit iostream_init;  // in each TU that includes the header
\end{lstlisting}

\subsection{Named Parameter Idiom}
\label{subsec:named-parameter-idiom}
See~\cite{isocpp-named-parameter-idiom}.
\newline
Emulate named/defaulted arguments by returning \texttt{*this} from chainable 
setters (a lightweight builder). Call sites read as a sequence of named 
options instead of a long positional argument list.

\begin{lstlisting}[language=C++]
class Game {
public:
    Game& windowSize(int w, int h) { w_ = w; h_ = h; return *this; }
    Game& fullscreen(bool v) { full_ = v; return *this; }
    Game& title(std::string t) { title_ = std::move(t); return *this; }

private:
    int w_ = 800, h_ = 600;
    bool full_ = false;
    std::string title_ = "game";
};

Game g;
g.windowSize(1920, 1080).fullscreen(true).title("Demo");
\end{lstlisting}

\subsection{Virtual Friend Function Idiom}
\label{subsec:virtual-friend-function}
See~\cite{isocpp-virtual-friend-function}.
\newline
Friends are not virtual. To get polymorphic behavior for operators like 
\texttt{operator<<}, provide a public non-virtual friend that forwards to 
a private/protected virtual member.

\begin{lstlisting}[language=C++]
class Shape {
public:
    virtual ~Shape() = default;
    friend std::ostream& operator<<(std::ostream& os, const Shape& s) {
        return s.print(os);          // virtual dispatch
    }
protected:
    virtual std::ostream& print(std::ostream& os) const = 0;
};

class Circle : public Shape {
protected:
    std::ostream& print(std::ostream& os) const override {
        return os << "Circle";
    }
};

void f(const Shape& s) { std::cout << s << '\n'; }
\end{lstlisting}

\subsection{CRTP - Curiously recurring template pattern}
\label{subsec:crtp}
See~\cite{wiki-crtp},~\cite{cppreference-crtp}.
\newline
A class \texttt{Derived} inherits from \texttt{Base<Derived>}. The base 
can \texttt{static\char`_cast} to \texttt{Derived} and call derived members 
\textbf{without virtual functions} (static polymorphism). Common uses: 
mixins, compile-time interfaces, \texttt{enable\char`_shared\char`_from\char`_this}-style 
helpers.

\begin{lstlisting}[language=C++]
template <typename D>
struct Comparable {
    bool operator==(const D& other) const {
        return static_cast<const D*>(this)->value()
            == other.value();
    }
    bool operator!=(const D& other) const { return !(*this == other); }
};

struct Point : Comparable<Point> {
    int x, y;
    int value() const { return x ^ y; }  // used by Comparable
};

Point a{1, 2}, b{1, 2};
bool same = (a == b);   // resolved at compile time
\end{lstlisting}

\subsection{Type Erasure}
\label{subsec:type-erasure}
See~\cite{wikibook-type-erasure}.
\newline
Hide a concrete type behind a uniform interface while retaining behaviour 
(virtuals, function pointers, or a small-buffer-optimized wrapper). Clients 
depend on the erased handle; the real type is known only at the construction 
site. \texttt{std::function}, \texttt{std::any}, and polymorphic allocators 
are standard examples.

\begin{lstlisting}[language=C++]
class Drawable {                 // type-erased handle
public:
    template <class T>
    Drawable(T x) : self_{std::make_shared<Model<T>>(std::move(x))} {}

    void draw() const { self_->draw(); }

private:
    struct Concept {
        virtual ~Concept() = default;
        virtual void draw() const = 0;
    };
    template <class T>
    struct Model : Concept {
        explicit Model(T x) : data_{std::move(x)} {}
        void draw() const override { data_.draw(); }
        T data_;
    };
    std::shared_ptr<const Concept> self_;
};

struct Circle { void draw() const { /* ... */ } };
struct Square { void draw() const { /* ... */ } };

std::vector<Drawable> scene{Circle{}, Square{}};
for (const auto& d : scene) d.draw();
\end{lstlisting}
